\documentclass[11pt]{article}
% Generated by scripts/build_checklists.py from the data files in data/. Do not edit by hand.
\pdfinfoomitdate=1
\pdftrailerid{}
\pdfsuppressptexinfo=-1
\usepackage[letterpaper,margin=0.8in]{geometry}
\usepackage[T1]{fontenc}
\usepackage{lmodern}
\usepackage{amssymb}
\usepackage{enumitem}
\usepackage[most]{tcolorbox}
\usepackage{xcolor}
\usepackage{titlesec}
\usepackage[hidelinks,pdfcreator={precalc-tutor},pdfproducer={pdflatex}]{hyperref}

\definecolor{accent}{HTML}{1F3A5F}
\definecolor{soft}{HTML}{EEF3FA}
\definecolor{warn}{HTML}{FFF4E5}
\pagestyle{plain}
\setlength{\parindent}{0pt}
\setlength{\parskip}{4pt}
\raggedright

\titleformat{\section}{\large\bfseries\color{accent}}{}{0pt}{}
\titlespacing*{\section}{0pt}{14pt}{4pt}
\titleformat{\subsection}{\normalsize\bfseries\color{accent}}{}{0pt}{}
\titlespacing*{\subsection}{0pt}{10pt}{2pt}

\newlist{checks}{itemize}{1}
\setlist[checks]{label=$\square$,leftmargin=2em,itemsep=2pt,topsep=2pt}
\newcommand{\aside}[1]{{\small\itshape #1}}
\newcommand{\tagright}[1]{\hfill{\small\color{gray}#1}}
\newtcolorbox{infobox}{colback=soft,colframe=accent,boxrule=0.6pt,arc=2pt,left=6pt,right=6pt,top=4pt,bottom=4pt}
\newtcolorbox{draftbox}{colback=warn,colframe=orange!70!black,boxrule=0.5pt,arc=2pt,left=6pt,right=6pt,top=3pt,bottom=3pt}

\hypersetup{pdftitle={Chapter 4 checklist: Trigonometry}}
\begin{document}

{\LARGE\bfseries\color{accent} Chapter 4 Checklist}\par
\vspace{2pt}
{\large Trigonometry}\par
\vspace{2pt}
{\small Ron Larson and Robert P. Hostetler, \emph{Precalculus with Limits}, 2nd ed., p. 279. Tick each box as you go.}\par
\vspace{6pt}
\begin{infobox}
Work through the four parts in order: read each section, practise the matching Duolingo units, work the book-only topics that Duolingo does not cover, then finish the chapter. Duolingo units are grouped by lane (a grade or a topic in the app) and listed in the order the app shows them, so you can follow one lane from top to bottom. \textbf{full} = the unit practises the section's main skill, \textbf{partial} = part of it, \textbf{prereq} = a warm-up from an earlier grade.
\end{infobox}
\textbf{In this chapter:} If time is short, follow grade 10 for the right-triangle ratios and special-angle values of 4.3, then grade 11 into grade 12 for the sine, cosine, and tangent graphs of 4.5 and 4.6, where Duolingo is strongest. No unit covers radians, the unit circle, reference angles, cosecant, secant, cotangent, or inverse trig functions, so most of 4.1, 4.2, 4.4, and 4.7 comes from the book, with Duolingo giving only single steps and warm-ups. Section 4.8 gets the tools but not the applications: elevation and depression problems, bearings, and harmonic motion models are book only.

\begin{draftbox}
\textbf{Draft.} The Duolingo mapping for this chapter has not been reviewed yet, so treat the unit choices as a first pass.
\end{draftbox}

\section{1.\quad Read the textbook sections}
\begin{checks}
  \item 4.1 Radian and Degree Measure\tagright{p. 280}
  \item 4.2 Trigonometric Functions: The Unit Circle\tagright{p. 292}
  \item 4.3 Right Triangle Trigonometry\tagright{p. 299}
  \item 4.4 Trigonometric Functions of Any Angle\tagright{p. 310}
  \item 4.5 Graphs of Sine and Cosine Functions\tagright{p. 319}
  \item 4.6 Graphs of Other Trigonometric Functions\tagright{p. 330}
  \item 4.7 Inverse Trigonometric Functions\tagright{p. 341}
  \item 4.8 Applications and Models\tagright{p. 351}
\end{checks}

\section{2.\quad Practise on Duolingo}
Pick one lane and work down it; the grade lanes go deepest. The note after each unit gives its rating and the textbook sections it serves.
\subsection{Grade 4}
\begin{checks}
  \item Time\quad\aside{prereq, 4.1}\tagright{unit 46}
  \item Angle measures\quad\aside{prereq, 4.1}\tagright{unit 51}
\end{checks}
\subsection{Grade 5}
\begin{checks}
  \item Add fractions with unlike denominators\quad\aside{prereq, 4.1}\tagright{unit 13}
  \item Subtract fractions with unlike denominators\quad\aside{prereq, 4.1}\tagright{unit 14}
  \item Multiply fractions\quad\aside{prereq, 4.1}\tagright{unit 21}
\end{checks}
\subsection{Grade 6}
\begin{checks}
  \item Unit rate\quad\aside{prereq, 4.1}\tagright{unit 10}
  \item Unit conversions in the customary system\quad\aside{prereq, 4.1}\tagright{unit 14}
  \item Greatest common factor\quad\aside{prereq, 4.1}\tagright{unit 30}
  \item Plot points in all quadrants\quad\aside{prereq, 4.1, 4.4}\tagright{unit 33}
\end{checks}
\subsection{Grade 7}
\begin{checks}
  \item Circumference of circles\quad\aside{prereq, 4.1, 4.2}\tagright{unit 41}
  \item Area of circles\quad\aside{prereq, 4.1}\tagright{unit 42}
  \item Complementary angles\quad\aside{partial, 4.1, 4.3, 4.8}\tagright{unit 43}
  \item Supplementary angles\quad\aside{partial, 4.1, 4.4}\tagright{unit 44}
\end{checks}
\subsection{Grade 8}
\begin{checks}
  \item Alternate interior angles\quad\aside{prereq, 4.8}\tagright{unit 73}
  \item Interior angles of a triangle\quad\aside{prereq, 4.8}\tagright{unit 76}
\end{checks}
\subsection{Grade 9}
\begin{checks}
  \item Simplify radicals\quad\aside{prereq, 4.3, 4.4}\tagright{unit 43}
\end{checks}
\subsection{Grade 10}
\begin{checks}
  \item Define rotations\quad\aside{prereq, 4.1}\tagright{unit 6}
  \item 90-degree clockwise rotations\quad\aside{prereq, 4.1}\tagright{unit 8}
  \item Proportional relationships in similar triangles\quad\aside{prereq, 4.3}\tagright{unit 17}
  \item Pythagorean theorem\quad\aside{prereq, 4.3, 4.4, 4.7, 4.8}\tagright{unit 20}
  \item Intro to the sine ratio\quad\aside{full, 4.3, 4.7, 4.8}\tagright{unit 23}
  \item Sine of 30\quad\aside{partial, 4.2, 4.3, 4.4}\tagright{unit 24}
  \item Sine of 60\quad\aside{partial, 4.2, 4.3, 4.4}\tagright{unit 25}
  \item Sine of 45\quad\aside{partial, 4.2, 4.3, 4.4}\tagright{unit 26}
  \item Intro to the cosine ratio\quad\aside{full, 4.3, 4.7, 4.8}\tagright{unit 27}
  \item Cosine of 60\quad\aside{partial, 4.2, 4.3, 4.4}\tagright{unit 28}
  \item Cosine of 30\quad\aside{partial, 4.2, 4.3, 4.4}\tagright{unit 29}
  \item Cosine of 45\quad\aside{partial, 4.2, 4.3, 4.4}\tagright{unit 30}
  \item Intro to the tangent ratio\quad\aside{full, 4.3, 4.7, 4.8}\tagright{unit 31}
  \item Tangent of 45\quad\aside{partial, 4.2, 4.3, 4.4}\tagright{unit 32}
  \item Tangent of 30\quad\aside{partial, 4.2, 4.3, 4.4}\tagright{unit 33}
  \item Tangent of 60\quad\aside{partial, 4.2, 4.3, 4.4}\tagright{unit 34}
  \item Connect sine and cosine\quad\aside{partial, 4.3, 4.7}\tagright{unit 35}
  \item Connect sine, cosine, and tangent\quad\aside{partial, 4.2, 4.3, 4.6, 4.7}\tagright{unit 36}
  \item Reflections across the x-axis as functions\quad\aside{prereq, 4.2}\tagright{unit 41}
  \item Reflections across y=x as functions\quad\aside{prereq, 4.7}\tagright{unit 43}
  \item Circles on the plane\quad\aside{prereq, 4.2}\tagright{unit 47}
  \item Distance formula\quad\aside{prereq, 4.4}\tagright{unit 51}
\end{checks}
\subsection{Grade 11}
\begin{checks}
  \item Exponential decay\quad\aside{prereq, 4.6}\tagright{unit 42}
  \item Horizontal translations\quad\aside{prereq, 4.5}\tagright{unit 70}
  \item X-axis function reflections\quad\aside{partial, 4.5, 4.6}\tagright{unit 71}
  \item Y-axis function reflections\quad\aside{partial, 4.5}\tagright{unit 72}
  \item Intro to the sine function\quad\aside{partial, 4.2, 4.5, 4.6, 4.7}\tagright{unit 77}
  \item Period of sine functions\quad\aside{partial, 4.2, 4.5}\tagright{unit 78}
  \item Intro to the cosine function\quad\aside{partial, 4.2, 4.5, 4.6, 4.7}\tagright{unit 79}
  \item Period of cosine functions\quad\aside{partial, 4.2, 4.5}\tagright{unit 80}
  \item Midline of sine and cosine functions\quad\aside{full, 4.5}\tagright{unit 81}
  \item Amplitude of sine and cosine functions\quad\aside{full, 4.5}\tagright{unit 82}
  \item Intro to the tangent function\quad\aside{partial, 4.2, 4.6, 4.7}\tagright{unit 83}
  \item Period of tangent functions\quad\aside{full, 4.6}\tagright{unit 84}
\end{checks}
\subsection{Grade 12}
\begin{checks}
  \item Vertical asymptotes of rational functions\quad\aside{prereq, 4.6}\tagright{unit 9}
  \item Function composition\quad\aside{partial, 4.7}\tagright{unit 13}
  \item Inverse functions\quad\aside{partial, 4.7}\tagright{unit 14}
  \item Domain and range\quad\aside{partial, 4.7}\tagright{unit 17}
  \item Graphing inverses\quad\aside{partial, 4.7}\tagright{unit 18}
  \item Period of sinusoidal functions\quad\aside{full, 4.5, 4.6, 4.8}\tagright{unit 19}
  \item Midline of sinusoidal functions\quad\aside{full, 4.5}\tagright{unit 20}
  \item Amplitude of sinusoidal functions\quad\aside{full, 4.5, 4.6, 4.8}\tagright{unit 21}
  \item Phase shift of sinusoidal functions\quad\aside{full, 4.5, 4.6}\tagright{unit 22}
  \item Relate sine and cosine\quad\aside{partial, 4.5, 4.8}\tagright{unit 23}
\end{checks}
\subsection{Geometry}
\begin{checks}
  \item Circles\quad\aside{prereq, 4.1}\tagright{unit 12}
  \item Angle relationships\quad\aside{partial, 4.1}\tagright{unit 13}
  \item The Pythagorean theorem\quad\aside{prereq, 4.3, 4.8}\tagright{unit 23}
  \item Sine ratio\quad\aside{full, 4.3, 4.8}\tagright{unit 24}
  \item Cosine ratio\quad\aside{full, 4.3, 4.8}\tagright{unit 25}
  \item Tangent ratio\quad\aside{full, 4.3, 4.8}\tagright{unit 26}
  \item Sine function\quad\aside{partial, 4.2, 4.5}\tagright{unit 27}
  \item Cosine function\quad\aside{partial, 4.2, 4.5}\tagright{unit 28}
  \item Tangent function\quad\aside{partial, 4.2, 4.6}\tagright{unit 29}
  \item Graph sine and cosine\quad\aside{full, 4.5, 4.6, 4.8}\tagright{unit 30}
\end{checks}
\subsection{Algebraic Graphing}
\begin{checks}
  \item Inverse linear functions\quad\aside{prereq, 4.7}\tagright{unit 23}
\end{checks}

\section{3.\quad Section by section}
For each section: what Duolingo covers (the units are ticked off in part 2) and the book-only topics to work in the textbook. Tick each book-only topic once you have worked examples of it.
\subsection{4.1 Radian and Degree Measure\tagright{p. 280}}
Duolingo has no unit on standard position, coterminal angles, radians, arc length, angular speed or sectors, so most of this section is book only. The grade 7 complementary and supplementary units practise that skill in degrees, with the Geometry topic's Angle relationships as the Topics-tab alternative. The angle, rotation, quadrant, fraction, factor, circle, conversion, time and rate units are warm-ups for the section's definitions and formulas.\par
{\small \textit{Duolingo:} \textbf{Grade 4:} 46. Time (prereq); 51. Angle measures (prereq)\\ \textbf{Grade 5:} 13. Add fractions with unlike denominators (prereq); 14. Subtract fractions with unlike denominators (prereq); 21. Multiply fractions (prereq)\\ \textbf{Grade 6:} 10. Unit rate (prereq); 14. Unit conversions in the customary system (prereq); 30. Greatest common factor (prereq); 33. Plot points in all quadrants (prereq)\\ \textbf{Grade 7:} 41. Circumference of circles (prereq); 42. Area of circles (prereq); 43. Complementary angles (partial); 44. Supplementary angles (partial)\\ \textbf{Grade 10:} 6. Define rotations (prereq); 8. 90-degree clockwise rotations (prereq)\\ \textbf{Geometry:} 12. Circles (prereq); 13. Angle relationships (partial)}\par
{\small\textit{Book only:}}
\begin{checks}
  \item Angles in standard position: vertex, initial side, terminal side
  \item Positive (counterclockwise) and negative (clockwise) angles
  \item Angles beyond one full revolution
  \item Quadrant of an angle's terminal side
  \item Coterminal angles in degrees and radians
  \item Definition of radian measure (arc length over radius)
  \item Radian measures of common angles and quadrant boundaries
  \item Acute and obtuse angles in radians (0 to pi/2, pi/2 to pi)
  \item Converting between degrees and radians
  \item Degrees-minutes-seconds and decimal degrees
  \item Complementary and supplementary angles in radians
  \item Arc length formula s = r theta
  \item Linear speed v = s/t, angular speed omega = theta/t, and v = r omega
  \item Converting revolutions per minute to radians per minute
  \item Area of a circular sector
\end{checks}
\subsection{4.2 Trigonometric Functions: The Unit Circle\tagright{p. 292}}
No unit is named for radians, the unit circle, or cosecant, secant, and cotangent, so much of this section is book only. Drill the exact first-quadrant values with the Grade 10 special-angle units, which work in degrees (match 30, 45, 60 to pi/6, pi/4, pi/3 yourself), and use Grade 11 for sine, cosine, and tangent as functions and for reading the period; the Geometry topic's three function units are an alternative for the first part only. Book only, among others listed in the gaps: defining all six functions from (x, y), values outside the first quadrant, even and odd functions, using the period to evaluate, and calculator work.\par
{\small \textit{Duolingo:} \textbf{Grade 7:} 41. Circumference of circles (prereq)\\ \textbf{Grade 10:} 24. Sine of 30 (partial); 25. Sine of 60 (partial); 26. Sine of 45 (partial); 28. Cosine of 60 (partial); 29. Cosine of 30 (partial); 30. Cosine of 45 (partial); 32. Tangent of 45 (partial); 33. Tangent of 30 (partial); 34. Tangent of 60 (partial); 36. Connect sine, cosine, and tangent (partial); 41. Reflections across the x-axis as functions (prereq); 47. Circles on the plane (prereq)\\ \textbf{Grade 11:} 77. Intro to the sine function (partial); 78. Period of sine functions (partial); 79. Intro to the cosine function (partial); 80. Period of cosine functions (partial); 83. Intro to the tangent function (partial)\\ \textbf{Geometry:} 27. Sine function (partial); 28. Cosine function (partial); 29. Tangent function (partial)}\par
{\small\textit{Book only:}}
\begin{checks}
  \item The unit circle and wrapping the real number line around it
  \item Defining the six trigonometric functions from the point (x, y) for a real number t
  \item Cosecant, secant, and cotangent values as reciprocals
  \item Special-angle values in radians (pi/6, pi/4, pi/3), matched to the degree drills
  \item Values at the quadrantal points 0, pi/2, pi, 3pi/2
  \item Values and signs at standard points outside the first quadrant
  \item Where tangent, cotangent, secant, and cosecant are undefined
  \item Domain (all reals) and range [-1, 1] of sine and cosine
  \item Using the 2pi period to evaluate at large or negative t
  \item Even and odd trigonometric functions (cosine and secant even, the others odd)
  \item Calculator evaluation of trig functions in radian mode, including csc, sec, cot as reciprocals
  \item Estimating sine and cosine values from a unit circle figure, and approximating t for a given value
  \item Evaluating trigonometric models at given times, such as a spring's displacement
\end{checks}
\subsection{4.3 Right Triangle Trigonometry\tagright{p. 299}}
Grade 10 covers the core in degrees: sine, cosine and tangent as right-triangle ratios, their exact values at 30, 45 and 60 degrees, sine and cosine of complementary angles, and tangent as sine over cosine; the Geometry topic tab is a shorter alternative for the three ratios. Cosecant, secant and cotangent, most identity work, radian forms of the special angles, calculator work, solving right triangles, and angle of elevation and depression problems are book only.\par
{\small \textit{Duolingo:} \textbf{Grade 7:} 43. Complementary angles (prereq)\\ \textbf{Grade 9:} 43. Simplify radicals (prereq)\\ \textbf{Grade 10:} 17. Proportional relationships in similar triangles (prereq); 20. Pythagorean theorem (prereq); 23. Intro to the sine ratio (full); 24. Sine of 30 (full); 25. Sine of 60 (full); 26. Sine of 45 (full); 27. Intro to the cosine ratio (full); 28. Cosine of 60 (full); 29. Cosine of 30 (full); 30. Cosine of 45 (full); 31. Intro to the tangent ratio (full); 32. Tangent of 45 (full); 33. Tangent of 30 (full); 34. Tangent of 60 (full); 35. Connect sine and cosine (partial); 36. Connect sine, cosine, and tangent (partial)\\ \textbf{Geometry:} 23. The Pythagorean theorem (prereq); 24. Sine ratio (full); 25. Cosine ratio (full); 26. Tangent ratio (full)}\par
{\small\textit{Book only:}}
\begin{checks}
  \item Cosecant, secant, and cotangent as right-triangle ratios
  \item Reciprocal identities
  \item Quotient identity for cotangent (cot = cos/sin)
  \item Pythagorean identities
  \item Cofunction identities beyond the sine-cosine pair (tangent-cotangent, secant-cosecant)
  \item Special-angle values in radians (pi/6, pi/4, pi/3), and finding the angle from a value
  \item Finding all six function values from one given ratio
  \item Using fundamental identities to simplify or verify trig expressions
  \item Calculator evaluation of trig functions at non-special acute angles (degree mode, decimal degrees and DMS, reciprocals via 1/x)
  \item Finding an acute angle from a known side ratio with a calculator
  \item Solving right triangles at non-special angles (missing sides)
  \item Angles of elevation and depression in applied height and distance problems
\end{checks}
\subsection{4.4 Trigonometric Functions of Any Angle\tagright{p. 310}}
Duolingo has no unit on angles in standard position, quadrant signs, or reference angles, so the section's method is book only. It supports single steps: the distance formula, Pythagorean theorem, and radical simplification give r or a missing coordinate, supplementary angles give a second-quadrant reference angle, and the grade 10 special-angle units give the value at a 30, 45, or 60 degree reference angle.\par
{\small \textit{Duolingo:} \textbf{Grade 6:} 33. Plot points in all quadrants (prereq)\\ \textbf{Grade 7:} 44. Supplementary angles (prereq)\\ \textbf{Grade 9:} 43. Simplify radicals (prereq)\\ \textbf{Grade 10:} 20. Pythagorean theorem (prereq); 24. Sine of 30 (partial); 25. Sine of 60 (partial); 26. Sine of 45 (partial); 28. Cosine of 60 (partial); 29. Cosine of 30 (partial); 30. Cosine of 45 (partial); 32. Tangent of 45 (partial); 33. Tangent of 30 (partial); 34. Tangent of 60 (partial); 51. Distance formula (prereq)}\par
{\small\textit{Book only:}}
\begin{checks}
  \item Six trig functions defined from a point (x, y) and r
  \item Cosecant, secant, and cotangent of any angle
  \item Signs of the trig functions by quadrant
  \item Identifying the quadrant from the signs of two functions
  \item Choosing a point on a terminal side given as a line in a stated quadrant
  \item Values at quadrantal angles, including where a function is undefined
  \item Finding the remaining function values from one value and its quadrant
  \item Finding reference angles in degrees and radians
  \item Coterminal step for negative angles and angles past one revolution
  \item Combining a reference-angle value with the quadrant sign
  \item Special-angle values in radians (pi/6, pi/4, pi/3)
  \item Finding both angles in one revolution with a given special value
  \item Trig functions of real numbers evaluated with reference angles
  \item Calculator evaluation of trig functions in degree and radian mode, including csc, sec, cot as reciprocals
  \item Evaluating trig models in applied problems
\end{checks}
\subsection{4.5 Graphs of Sine and Cosine Functions\tagright{p. 319}}
Grade 11 introduces the sine and cosine graphs with period, midline, and amplitude, and grade 12 repeats them for full sinusoids and adds phase shift, so do grade 11 first, or use the Geometry topic's three sine and cosine units as a shorter alternative. No unit reflects a sinusoid by name, so the general grade 11 reflection units stand in, and the grade 11 horizontal translation unit is a warm-up for phase shift. Key-point sketching, rewriting a negative b with the even/odd properties, frequency, writing an equation from a graph, and fitting sinusoidal models to data are book only.\par
{\small \textit{Duolingo:} \textbf{Grade 11:} 70. Horizontal translations (prereq); 71. X-axis function reflections (partial); 72. Y-axis function reflections (partial); 77. Intro to the sine function (full); 78. Period of sine functions (full); 79. Intro to the cosine function (full); 80. Period of cosine functions (full); 81. Midline of sine and cosine functions (full); 82. Amplitude of sine and cosine functions (full)\\ \textbf{Grade 12:} 19. Period of sinusoidal functions (full); 20. Midline of sinusoidal functions (full); 21. Amplitude of sinusoidal functions (full); 22. Phase shift of sinusoidal functions (full); 23. Relate sine and cosine (partial)\\ \textbf{Geometry:} 27. Sine function (full); 28. Cosine function (full); 30. Graph sine and cosine (full)}\par
{\small\textit{Book only:}}
\begin{checks}
  \item Hand sketching a transformed sinusoid: finding the one-cycle interval and the five key points at quarter-period spacing
  \item Rewriting a negative b with the even/odd properties of sine and cosine
  \item Frequency as the reciprocal of the period
  \item Writing a complete sinusoid equation from its graph (choosing sine or cosine form and combining a, b, c, d)
  \item Sinusoidal models fitted to real data (finding a, b, c, d from a table or a description)
\end{checks}
\subsection{4.6 Graphs of Other Trigonometric Functions\tagright{p. 330}}
Grade 11 covers the parent tangent graph and its period; prefer that pair over the Topics-tab Geometry unit. The other units are warm-ups for guide curves, transformations, asymptotes and damping factors. Shifts, stretches, and asymptote locations of transformed tangent graphs, cotangent, secant, cosecant, damped graphs and the six-function summary are book only, and grade 12 "Secant lines" is calculus, not the secant function.\par
{\small \textit{Duolingo:} \textbf{Grade 10:} 36. Connect sine, cosine, and tangent (prereq)\\ \textbf{Grade 11:} 42. Exponential decay (prereq); 71. X-axis function reflections (partial); 77. Intro to the sine function (prereq); 79. Intro to the cosine function (prereq); 83. Intro to the tangent function (full); 84. Period of tangent functions (full)\\ \textbf{Grade 12:} 9. Vertical asymptotes of rational functions (prereq); 19. Period of sinusoidal functions (prereq); 21. Amplitude of sinusoidal functions (prereq); 22. Phase shift of sinusoidal functions (prereq)\\ \textbf{Geometry:} 29. Tangent function (full); 30. Graph sine and cosine (prereq)}\par
{\small\textit{Book only:}}
\begin{checks}
  \item Cotangent graph, period, and asymptotes
  \item Secant and cosecant graphs sketched from cosine and sine guide curves
  \item Asymptotes and local maxima and minima of secant and cosecant branches
  \item Vertical stretch of tangent and cotangent graphs, which have no amplitude
  \item Locating two consecutive asymptotes of transformed tangent and cotangent graphs by solving bx - c = -pi/2 and pi/2, or 0 and pi
  \item Phase shifts of tangent, cotangent, secant, and cosecant graphs
  \item Even/odd symmetry seen in the tangent, cotangent, secant, and cosecant graphs
  \item Domain, range, and period summary of all six trigonometric functions
  \item Solving equations such as tan x = c or sec x = c graphically on an interval
  \item Damped trigonometric graphs and damping factors
  \item Bounding curves y = plus or minus f(x) for products such as x sin x
  \item Behavior of functions such as (sin x)/x as x approaches 0
\end{checks}
\subsection{4.7 Inverse Trigonometric Functions\tagright{p. 341}}
No Duolingo unit teaches arcsin, arccos, or arctan, so the inverse trig functions themselves are book only. Prefer the grade 12 inverse-function block for the general mechanics, with the grade 10 right-triangle ratios and grade 11 trig graphs as warm-ups and the Topics-tab Inverse linear functions unit as a lighter linear-only start. Reading the grade 10 special-angle units from 4.3 (Sine of 30 and so on) backwards gives only acute angles in degrees for positive inputs; radian answers, negative inputs, and choosing the angle inside the restricted range are book only.\par
{\small \textit{Duolingo:} \textbf{Grade 10:} 20. Pythagorean theorem (prereq); 23. Intro to the sine ratio (prereq); 27. Intro to the cosine ratio (prereq); 31. Intro to the tangent ratio (prereq); 35. Connect sine and cosine (prereq); 36. Connect sine, cosine, and tangent (prereq); 43. Reflections across y=x as functions (prereq)\\ \textbf{Grade 11:} 77. Intro to the sine function (prereq); 79. Intro to the cosine function (prereq); 83. Intro to the tangent function (prereq)\\ \textbf{Grade 12:} 13. Function composition (partial); 14. Inverse functions (partial); 17. Domain and range (partial); 18. Graphing inverses (partial)\\ \textbf{Algebraic Graphing:} 23. Inverse linear functions (prereq)}\par
{\small\textit{Book only:}}
\begin{checks}
  \item Restricting sine, cosine, and tangent to one-to-one intervals
  \item Definitions, domains, and ranges of arcsin, arccos, and arctan
  \item The notation sin\textasciicircum{}-1 x and why it is not the reciprocal 1/sin x
  \item Exact values of arcsin, arccos, and arctan in radians, including negative and quadrantal inputs
  \item Calculator evaluation of inverse trig functions
  \item Graphs of y = arcsin x, y = arccos x, and y = arctan x, including transformed versions
  \item Inverse properties such as sin(arcsin x) = x and arcsin(sin y) = y, with their interval conditions
  \item Compositions whose input lies outside the restricted domain or range
  \item Exact values of compositions like tan(arccos a) at a number a, from a right-triangle sketch, including negative a
  \item Rewriting compositions such as sin(arctan x) as algebraic expressions in x
  \item Writing a triangle angle as an inverse trig function of x
  \item Rewriting one inverse trig expression as another via a right triangle (exercise extension)
  \item Inverse cotangent, secant, and cosecant (exercise extension)
  \item Identities such as arcsin(-x) = -arcsin x and arcsin x + arccos x = pi/2
  \item Applied problems that find an angle with an inverse trig function
\end{checks}
\subsection{4.8 Applications and Models\tagright{p. 351}}
Duolingo gives the tools, not the applications: the grade 10 right-triangle ratios (special angles only) support solving right triangles, and the grade 12 period, amplitude, and relate-sine-and-cosine units support reading harmonic motion models. Prefer the grade 10 and grade 12 lanes; the Geometry topic (The Pythagorean theorem through Tangent ratio, then Graph sine and cosine) is the Topics-tab alternative for both halves. Calculator and inverse-trig solving, elevation and depression problems, bearings, and harmonic motion itself (writing and evaluating models, frequency) are book only.\par
{\small \textit{Duolingo:} \textbf{Grade 7:} 43. Complementary angles (prereq)\\ \textbf{Grade 8:} 73. Alternate interior angles (prereq); 76. Interior angles of a triangle (prereq)\\ \textbf{Grade 10:} 20. Pythagorean theorem (prereq); 23. Intro to the sine ratio (partial); 27. Intro to the cosine ratio (partial); 31. Intro to the tangent ratio (partial)\\ \textbf{Grade 12:} 19. Period of sinusoidal functions (partial); 21. Amplitude of sinusoidal functions (partial); 23. Relate sine and cosine (partial)\\ \textbf{Geometry:} 23. The Pythagorean theorem (prereq); 24. Sine ratio (partial); 25. Cosine ratio (partial); 26. Tangent ratio (partial); 30. Graph sine and cosine (partial)}\par
{\small\textit{Book only:}}
\begin{checks}
  \item Solving right triangles at non-special angles with a calculator
  \item Finding acute angles with inverse trig functions
  \item Angles of elevation and depression in applied height and distance problems
  \item Two-triangle problems, such as two angles of elevation to the same object
  \item Bearings in navigation (N/S-E/W compass form and three-digit nautical form)
  \item Simple harmonic motion models d = a sin(wt) and d = a cos(wt)
  \item Frequency of harmonic motion (w / 2pi)
  \item Writing a harmonic motion model from displacement at t = 0, amplitude, and period (choosing sine or cosine)
  \item Evaluating a harmonic motion model at a given time, including the least positive time when d = 0
\end{checks}

\section{4.\quad Finish the chapter}
\begin{checks}
  \item Chapter Summary\tagright{p. 362}
  \item Review Exercises\tagright{p. 364}
  \item Chapter Test\tagright{p. 367}
  \item Proofs in Mathematics\tagright{p. 368}
  \item Problem Solving\tagright{p. 369}
\end{checks}

\vfill
{\footnotesize\color{gray}
Duolingo unit numbers are each unit's position in that grade's or topic's path in the Duolingo Math app, as recorded on 2026-10-07 (iOS); the app may reorder them. Textbook: Ron Larson and Robert P. Hostetler. Precalculus with Limits, 2nd ed. (Pre-AP National Student Edition) Thomson Brooks/Cole, 2009. ISBN 978-0-538-73659-6. Only titles and unit names are reproduced, for reference; no textbook or Duolingo content is included. This checklist is licensed CC BY 4.0; see the repository NOTICE for third-party material.\par}

\end{document}
